\section{Notación básica}
\subsection{Conjuntos}
Sea $\Gamma=\set{k_1,\dots,k_l}\subseteq\N_n$ un conjunto con cardinal $|\Gamma|=l$ y $\gamma=(\gamma_{k})_{k=1}^l\in\B^l$, denotamos por $\ext{\gamma}{\Gamma}=(\ext{\gamma}{\Gamma}_j)_{j=1}^n$ el elemento de $\mathbb{Z}_2^n$ definido por:
\begin{equation}
  \label{eq:vector_extendido}
  \ext{\gamma}{\Gamma}_j=
  \begin{cases}
    \gamma_l & \text{ si } j =k_l\in\Gamma \\
    0        & \text{ si }j\notin\Gamma
  \end{cases}\,.
\end{equation}

Para el proceso inverso, dado $\gamma=(\gamma_{j})_{j=1}^n\in\B^n$, llamaremos $\con{\gamma}{\Gamma}$ al elemento de $\B^l$ definido por
\begin{equation}
  \label{eq:vector_reducido}
  \con{\gamma}{\Gamma}=(\gamma_{k_m})_{m=1}^l\,.
\end{equation}


Por ejemplo, para $n=8$, si $\Gamma = \set{1, 3, 7}$ y $\Delta = \set{2, 4}$, dados $\gamma = (0,1,1)\in\B^3$ y $\delta = (1, 1, 1, 0, 0, 1, 1, 0)\in\B^8$, los elemento $\ext{\gamma}{\Gamma}$ y $\con{\delta}{\Delta}$ serán:
\begin{align*}
  \ext{\gamma}{\Gamma} & = (0, 0, 1, 0, 0, 0, 1, 0)\,, \\
  \con{\delta}{\Delta} & = (1, 0)\,.
\end{align*}


Observemos que para todo elemento $\gamma\in\mathbb{Z}^l$ se cumple que $\con{\ext{\gamma}{\Gamma}}{\Gamma}=\gamma$ mientras que en general si $\gamma\in\mathbb{Z}^n$ entonces $\ext{\con{\gamma}{\Gamma}}{\Gamma}\neq\gamma$.

Cuando un conjunto tenga solo un elemento, se pude usar dicho elemento como identificación del conjunto, y por tanto, sin que suponga un conflicto de interpretación, si $k\in\N$ es un número natural, usaremos este número para representar al conjunto $\set{k}$.

También vamos a considerar que los elementos $(0,\dots,0)\in\B^n$ y $(1,\dots,1)\in\B^n$ lo denotaremos simplemente por $0$ y $1$ y como $0_n$ y $1_n$ cuando sea importante indicar el número de coordenadas.
